\chapter{Especificación de requisitos}

\section{Actores}

El sistema tiene un único actor humano y se apoya en dos sistemas externos
con los que intercambia información a lo largo de una sesión.

\begin{itemize}
\item \textbf{Usuario investigador} (actor primario). Persona que describe
un problema de decisión en lenguaje natural, lanza simulaciones, sigue su
ejecución en tiempo real y consulta los análisis, los informes y el
historial de experimentos. Es el único actor que interactúa con la
interfaz web.
\item \textbf{Primera fase del proyecto} (sistema externo). Produce y
registra los modelos de decisión que el laboratorio descubre y carga
dinámicamente. No actúa durante la sesión, pero su salida condiciona qué
modelos están disponibles.
\item \textbf{Servicios de modelos de lenguaje y apoyo} (sistema externo).
Proveedores accedidos a través de OpenRouter para el razonamiento de los
agentes, y el servicio opcional de memoria semántica
(\textit{embeddings} con Voyage AI y \textit{reranking} con ZeroEntropy)
que da soporte a la recuperación avanzada de conocimiento y a la escritura
indexada de memorias.
\end{itemize}

\section{Casos de uso}

Los casos de uso recogen las interacciones que el usuario investigador puede
realizar desde el laboratorio. En todos ellos la entrada pasa por una sesión de
chat con el Orchestrator, que decide qué agentes y herramientas intervienen en
cada momento.

\begin{figure}[H]
\centering
\resizebox{0.96\linewidth}{!}{%
\begin{tikzpicture}[
  actorlabel/.style={font=\small\sffamily, align=center},
  usecase/.style={
    ellipse, draw=black!75, line width=0.55pt, fill=black!3,
    minimum width=36mm, minimum height=10mm,
    align=center, font=\small\sffamily
  },
  systembox/.style={
    rectangle, draw=black!60, rounded corners=2pt, line width=0.5pt,
    fill=black!1, minimum width=104mm, minimum height=72mm
  },
  ext/.style={
    rectangle, draw=black!65, rounded corners=2pt, line width=0.5pt,
    fill=black!4, text width=28mm, minimum height=9mm,
    align=center, font=\scriptsize\sffamily
  },
  assoc/.style={draw=black!65, line width=0.5pt},
  include/.style={-{Stealth[length=2mm,width=1.6mm]}, draw=black!55,
    dashed, line width=0.45pt}
]
  \node[systembox] (system) at (4.6,0) {};
  \node[font=\small\sffamily\bfseries, anchor=north west] at
    ([xshift=2mm,yshift=-2mm]system.north west) {Laboratorio virtual};

  \coordinate (head) at (-2.0,1.7);
  \draw[assoc] (head) circle (3mm);
  \draw[assoc] (-2.0,1.4) -- (-2.0,0.4);
  \draw[assoc] (-2.55,1.05) -- (-1.45,1.05);
  \draw[assoc] (-2.0,0.4) -- (-2.45,-0.35);
  \draw[assoc] (-2.0,0.4) -- (-1.55,-0.35);
  \node[actorlabel, below] at (-2.0,-0.45) {Usuario\\investigador};

  \node[usecase] (cu1) at (3.2,2.25) {Ejecutar\\simulación};
  \node[usecase] (cu2) at (6.1,2.25) {Comparar\\modelos};
  \node[usecase] (cu3) at (3.2,0.55) {Analizar\\resultados};
  \node[usecase] (cu4) at (6.1,0.55) {Generar\\informe PDF};
  \node[usecase] (cu5) at (3.2,-1.2) {Consultar\\experimentos};
  \node[usecase] (cu6) at (6.1,-1.2) {Inspeccionar\\Knowledge Backbone};

  \node[ext] (phase1) at (10.0,1.65) {Primera fase\\modelos aceptados};
  \node[ext] (llm) at (10.0,-0.15) {OpenRouter\\y servicios semánticos};

  \draw[assoc] (-1.4,1.1) -- (cu1.west);
  \draw[assoc] (-1.4,0.85) -- (cu2.west);
  \draw[assoc] (-1.4,0.55) -- (cu3.west);
  \draw[assoc] (-1.4,0.25) -- (cu4.west);
  \draw[assoc] (-1.4,-0.05) -- (cu5.west);
  \draw[assoc] (-1.4,-0.35) -- (cu6.west);

  \draw[include] (cu2) -- node[above, font=\scriptsize] {<<include>>} (cu1);
  \draw[include] (cu4) -- node[right, font=\scriptsize] {<<include>>} (cu3);
  \draw[assoc] (phase1) -- (cu1.east);
  \draw[assoc] (phase1) -- (cu2.east);
  \draw[assoc] (llm) -- (cu1.east);
  \draw[assoc] (llm) -- (cu3.east);
  \draw[assoc] (llm) -- (cu5.east);
\end{tikzpicture}}
\caption[Diagrama de casos de uso]{Diagrama de casos de uso UML de la
segunda fase. Los sistemas externos quedan fuera del límite del laboratorio
porque el usuario no los opera directamente.}
\label{fig:casos-uso}
\end{figure}

\begin{table}[H]
\begin{center}
\small
\setlength{\tabcolsep}{4pt}
\renewcommand{\arraystretch}{1.12}
\begin{tabular}{|>{\raggedright\arraybackslash}p{1.45cm}|>{\raggedright\arraybackslash}p{6.6cm}|>{\raggedright\arraybackslash}p{3.0cm}|} \hline
\textbf{Código} & \textbf{Caso de uso} & \textbf{RF relacionados} \\ \hline
CU-01 & Ejecutar una simulación a partir de una descripción en lenguaje natural & RF-01, RF-02, RF-03, RF-04, RF-10, RF-11, RF-12 \\ \hline
CU-02 & Comparar varios modelos en el mismo entorno & RF-01, RF-03, RF-06 \\ \hline
CU-03 & Analizar los resultados y contrastarlos con conocimiento previo & RF-04, RF-05, RF-06, RF-08 \\ \hline
CU-04 & Generar el informe PDF de la sesión & RF-06, RF-07 \\ \hline
CU-05 & Consultar y comparar experimentos anteriores & RF-09, RNF-11 \\ \hline
CU-06 & Inspeccionar el Knowledge Backbone desde la interfaz & RF-13 \\ \hline
\end{tabular}
\caption[Casos de uso]{Casos de uso y trazabilidad con requisitos
funcionales.}
\end{center}
\end{table}

CU-01 es el recorrido base: el usuario describe un problema, el Orchestrator
propone modelos y entorno, el Architect genera la especificación, se ejecuta la
simulación y la interfaz muestra los eventos en tiempo real. Los demás casos
de uso nacen de ese flujo. CU-02 compara modelos bajo la misma semilla y el
mismo entorno; CU-03 interpreta la simulación y escribe memoria reutilizable;
CU-04 convierte el análisis en un PDF; CU-05 recupera experimentos anteriores;
y CU-06 permite revisar procedencia, observaciones y relaciones del Knowledge
Backbone.

\section{Requisitos funcionales}

Los requisitos funcionales concretan qué debe poder hacer el laboratorio
durante una sesión completa: recibir la descripción del problema, preparar y
ejecutar la simulación, observar sus resultados, analizarlos y conservar la
evidencia para consultas posteriores. La prioridad se expresa con MoSCoW:
\emph{Must} para lo imprescindible, \emph{Should} para lo recomendable y
\emph{Could} para mejoras útiles pero no necesarias para cerrar el alcance.

\begin{table}[H]
\begin{center}
\scriptsize
\setlength{\tabcolsep}{3pt}
\renewcommand{\arraystretch}{1.14}
\begin{tabular}{|>{\raggedright\arraybackslash}p{1.1cm}|>{\raggedright\arraybackslash}p{2.8cm}|>{\raggedright\arraybackslash}p{6.9cm}|>{\raggedright\arraybackslash}p{1.25cm}|} \hline
\textbf{Código} & \textbf{Nombre} & \textbf{Descripción} & \textbf{Prio.} \\ \hline
RF-01 & Carga dinámica de modelos & Cargar los modelos de decisión generados por la primera fase, descubriéndolos sin intervención manual entre sesiones. & Must \\ \hline
RF-02 & Generación de entornos & Generar especificaciones de entorno de simulación a partir de descripciones en lenguaje natural. & Must \\ \hline
RF-03 & Ejecución de simulaciones & Ejecutar simulaciones con uno o varios modelos en el mismo entorno, soportando comparativas multi-modelo. & Must \\ \hline
RF-04 & Registro estructurado & Registrar eventos, episodios y trayectorias de decisión producidos durante cada simulación. & Must \\ \hline
RF-05 & Identificación de patrones & Identificar patrones de comportamiento, pasos críticos y eventos relevantes, contrastándolos con conocimiento previo. & Must \\ \hline
RF-06 & Visualizaciones analíticas & Generar gráficas y visualizaciones sobre trayectorias y métricas para incorporarlas al informe final. & Should \\ \hline
RF-07 & Generación de informes & Generar informes en PDF con resultados, evidencia visual y referencias a los paradigmas analizados. & Must \\ \hline
RF-08 & Persistencia de observaciones & Persistir observaciones de simulación en el Knowledge Backbone como memorias reutilizables. & Must \\ \hline
RF-09 & Historial de experimentos & Conservar el historial completo de experimentos y permitir su consulta, filtrado y comparativa, también mediante preguntas en lenguaje natural resolubles con SQL. & Should \\ \hline
RF-10 & Interacción conversacional & Permitir la interacción mediante un chat gestionado por el Orchestrator, con continuidad de conversación dentro de la sesión. & Must \\ \hline
RF-11 & Recomendaciones del Orchestrator & Recomendar al usuario qué hacer a continuación, qué modelos encajan y qué herramientas tiene disponibles. & Could \\ \hline
RF-12 & Replay de simulación & Ofrecer una visualización del entorno con \textit{replay} paso a paso y paneles de actividad por agente. & Should \\ \hline
RF-13 & Inspección del Knowledge Backbone & Exponer vistas de inspección del Knowledge Backbone, incluyendo grafo, memorias y procedencia, desde la interfaz web. & Could \\ \hline
\end{tabular}
\caption[Requisitos funcionales]{Requisitos funcionales de la segunda fase.}
\label{tab:requisitos-funcionales}
\end{center}
\end{table}

La tabla conecta los objetivos de la introducción con funciones verificables.
RF-02 y RF-12 cubren la definición y visualización del entorno; RF-01, RF-03 y
RF-04 cubren carga, ejecución y registro; RF-05 y RF-06 cubren el análisis;
RF-07 cierra el informe; y RF-08 y RF-09 mantienen la memoria reutilizable. El
uso diario del laboratorio queda recogido en RF-10, RF-11 y RF-13: conversación,
recomendaciones e inspección desde la interfaz. La trazabilidad con evidencias
de validación se recoge en el Apéndice~\ref{ap:trazabilidad-requisitos}.

\section{Requisitos no funcionales}

Los requisitos no funcionales fijan condiciones de calidad que atraviesan todo
el sistema. No describen una pantalla concreta ni una acción puntual del
usuario; describen cómo debe comportarse el laboratorio mientras carga modelos,
ejecuta simulaciones, conversa con el usuario y consulta servicios externos.

\begin{table}[H]
\begin{center}
\scriptsize
\setlength{\tabcolsep}{3pt}
\renewcommand{\arraystretch}{1.14}
\begin{tabular}{|>{\raggedright\arraybackslash}p{1.15cm}|>{\raggedright\arraybackslash}p{2.75cm}|>{\raggedright\arraybackslash}p{6.85cm}|>{\raggedright\arraybackslash}p{1.25cm}|} \hline
\textbf{Código} & \textbf{Nombre} & \textbf{Descripción} & \textbf{Prio.} \\ \hline
RNF-01 & Modularidad y desacoplamiento & Integrar la primera fase mediante un contrato mínimo de \textit{duck typing}, sin imponer dependencias de paquetes ni jerarquías de clases. & Must \\ \hline
RNF-02 & Extensibilidad de modelos & Incorporar nuevos modelos generados por la primera fase sin modificar el código de la segunda fase. & Must \\ \hline
RNF-03 & Trazabilidad & Conservar configuración, eventos y artefactos finales de cada ejecución con identificadores únicos auditables. & Must \\ \hline
RNF-04 & Reactividad & Enviar al cliente web actualizaciones de simulación y actividad de agentes en tiempo real mediante WebSocket, sin \textit{polling}. & Must \\ \hline
RNF-05 & Degradación controlada & Mantener operativo el flujo principal aunque no estén disponibles \textit{embeddings}, \textit{reranking} o escritura semántica en Qdrant. & Should \\ \hline
RNF-06 & Determinismo & Producir resultados idénticos para una misma semilla y configuración, facilitando comparativas reproducibles. & Must \\ \hline
RNF-07 & Portabilidad de despliegue & Levantar PostgreSQL, MinIO, Neo4j y Qdrant de forma reproducible mediante \texttt{docker compose}. & Should \\ \hline
RNF-08 & Configuración por entorno & Aislar credenciales y \textit{endpoints} en variables de entorno, sin embebidos en código ni artefactos versionados. & Must \\ \hline
RNF-09 & Observabilidad & Registrar acciones del Orchestrator y agentes especializados en los \textit{logs} del backend para diagnosticar sesiones \textit{a posteriori}. & Should \\ \hline
RNF-10 & Acotación del contexto conversacional & Mantener el historial del chat dentro de un presupuesto acotado automáticamente, sin reiniciar la sesión ni perder trazabilidad. & Should \\ \hline
RNF-11 & Recuperación relacional segura & Traducir preguntas en lenguaje natural a SQL de solo lectura, limitado a tablas permitidas y con máximo de resultados. & Must \\ \hline
RNF-12 & Contraste previo con predicciones & Recuperar predicciones científicas antes de ejecutar la simulación para compararlas después con la trayectoria observada. & Should \\ \hline
\end{tabular}
\caption[Requisitos no funcionales]{Requisitos no funcionales de la segunda
fase.}
\label{tab:requisitos-no-funcionales}
\end{center}
\end{table}

Estos requisitos protegen las decisiones principales del diseño. La segunda
fase no debe quedar atada al código interno de la primera; una simulación debe
poder auditarse después de ejecutarse; la interfaz debe recibir estado vivo por
WebSocket; y la ausencia de la capa semántica no debe impedir que el flujo
principal termine. También se acota el contexto conversacional, se limita la
consulta NL-SQL a lecturas seguras y se recuperan predicciones antes de observar
la simulación, para que el contraste posterior no dependa de una explicación
hecha a posteriori.
